\section{Reflections on GA2 and GA7}\label{sec:reflect}

\subsection{\gaTwoName}
Building the system meant applying what we had learned in programming and software design to a concrete problem. We used object-oriented decomposition to give every class a single responsibility, inheritance and abstraction to separate lecturer and Head of Department behaviour, and encapsulation to protect assessment weights. We reasoned numerically about the weighted mark and showed that scaling by the maximum mark is required, and we analysed the cost of each operation to justify choosing simple lists and linear passes over more complex structures. Working with \texttt{java.nio} and CSV forced us to deal with real data problems such as commas inside names, missing values and files whose columns are in a different order.

The most useful lesson was that applying knowledge is not the same as trusting it. The first version of the final-mark calculation looked correct and compiled, but it was wrong whenever assessments had different maximum marks. We only found it by comparing program output with hand-calculated values. After that we wrote the expected result before running each test, and we recorded our reasoning in the design decisions in Section~\ref{sec:decisions} so that it could be challenged.

\subsection{\gaSevenName}
Marks decide whether students progress, so an error in this software has consequences for people. We treated correctness as an ethical requirement as well as a technical one. The system refuses to calculate a final mark from incomplete data, rejects marks outside the valid range, and keeps missing marks out of class averages. Role-based access means the Head of Department can inspect results but cannot change them, which limits accidental or unauthorised edits.

Marks are also personal information, and in South Africa their handling falls under the Protection of Personal Information Act. In its current form the prototype stores plain-text CSV files with no authentication or encryption, and it keeps no audit trail of who changed a mark. Those are real limits, which is why the password feature is on its own branch for review and why we do not claim the system is ready for production use. On the positive side, the application needs only a standard Java installation and no server, the red highlighting and specific error messages make mistakes visible to non-technical users, and the modular design with explained code should keep long-term maintenance costs low.
